
\documentclass{article}
\usepackage{colortbl}
\usepackage{makecell}
\usepackage{multirow}
\usepackage{supertabular}

\begin{document}

\newcounter{utterance}

\centering \large Interaction Transcript for game `cladder', experiment `full\_v1.5\_default', episode 4315 with selfrefine{-}Qwen3.5{-}9B.
\vspace{24pt}

{ \footnotesize  \setcounter{utterance}{1}
\setlength{\tabcolsep}{0pt}
\begin{supertabular}{c@{$\;$}|p{.15\linewidth}@{}p{.15\linewidth}p{.15\linewidth}p{.15\linewidth}p{.15\linewidth}p{.15\linewidth}}
    \# & \multicolumn{2}{c}{Player} && \multicolumn{2}{c}{Game Master} \\
    \hline

    \theutterance \stepcounter{utterance}  
    & & & \multicolumn{4}{p{0.6\linewidth}}{
        \cellcolor[rgb]{0.9,0.9,0.9}{
            \makecell[{{p{\linewidth}}}]{
                \texttt{\tiny{[P1$\langle$GM]}}
                \texttt{You are an expert in causal inference. The following question is not a typical commonsense query, but rather a meticulously designed question created by a professor specializing in causal inference, intended to assess the students' mastery of the course content.} \\
\texttt{Imagine a self{-}contained, hypothetical world with only the following conditions, and without any unmentioned factors or causal relationships: Maternal smoking status has a direct effect on infant's birth weight and freckles. Health condition has a direct effect on infant's birth weight and freckles. Infant's birth weight has a direct effect on freckles. Health condition is unobserved. For infants with nonsmoking mothers and low infant birth weight, the probability of freckles is 46\%. For infants with nonsmoking mothers and normal infant birth weight, the probability of freckles is 21\%. For infants with smoking mothers and low infant birth weight, the probability of freckles is 74\%. For infants with smoking mothers and normal infant birth weight, the probability of freckles is 47\%. For infants with nonsmoking mothers and with poor health, the probability of normal infant birth weight is 51\%. For infants with nonsmoking mothers and with good health, the probability of normal infant birth weight is 72\%. For infants with smoking mothers and with poor health, the probability of normal infant birth weight is 3\%. For infants with smoking mothers and with good health, the probability of normal infant birth weight is 35\%. The overall probability of good health is 90\%. If we disregard the mediation effect through infant's birth weight, would maternal smoking status negatively affect freckles?} \\
\texttt{Start your answer with 'Yes' or 'No', followed by additional reasoning or evidence to support your explanation.} \\
            }
        }
    }
     \\ \\

    \theutterance \stepcounter{utterance}  
    & \multicolumn{4}{p{0.6\linewidth}}{
        \cellcolor[rgb]{0.561,0.737,0.561}{
            \makecell[{{p{\linewidth}}}]{
                \texttt{\tiny{[P1$\rangle$GM]}}
                \texttt{No} \\
\\ 
\texttt{**Reasoning:**} \\
\texttt{To determine if maternal smoking status negatively affects freckles while disregarding the mediation effect through infant's birth weight, we must calculate the **Direct Effect** of smoking on freckles, holding birth weight constant. This involves comparing the probability of freckles for smoking mothers versus nonsmoking mothers at the same levels of birth weight.} \\
\\ 
\texttt{From the provided data:} \\
\texttt{1.  **For infants with Low Birth Weight (\$W=0\$):**} \\
\texttt{    *   \$P(\textbackslash{}text\{Freckles\} | \textbackslash{}text\{Smoking\}, \textbackslash{}text\{Low Weight\}) = 74\textbackslash{}\%\$} \\
\texttt{    *   \$P(\textbackslash{}text\{Freckles\} | \textbackslash{}text\{Nonsmoking\}, \textbackslash{}text\{Low Weight\}) = 46\textbackslash{}\%\$} \\
\texttt{    *   Difference: \$74\textbackslash{}\% {-} 46\textbackslash{}\% = +28\textbackslash{}\%\$} \\
\\ 
\texttt{2.  **For infants with Normal Birth Weight (\$W=1\$):**} \\
\texttt{    *   \$P(\textbackslash{}text\{Freckles\} | \textbackslash{}text\{Smoking\}, \textbackslash{}text\{Normal Weight\}) = 47\textbackslash{}\%\$} \\
\texttt{    *   \$P(\textbackslash{}text\{Freckles\} | \textbackslash{}text\{Nonsmoking\}, \textbackslash{}text\{Normal Weight\}) = 21\textbackslash{}\%\$} \\
\texttt{    *   Difference: \$47\textbackslash{}\% {-} 21\textbackslash{}\% = +26\textbackslash{}\%\$} \\
\\ 
\texttt{In both strata of birth weight, the probability of having freckles is higher for infants of smoking mothers compared to nonsmoking mothers. Since the direct effect is positive in both cases, disregarding the mediation through birth weight reveals that maternal smoking status **positively** affects freckles, not negatively. Therefore, the answer to "would maternal smoking status negatively affect freckles?" is No.} \\
            }
        }
    }
    & & \\ \\

    \theutterance \stepcounter{utterance}  
    & & & \multicolumn{2}{p{0.3\linewidth}}{
        \cellcolor[rgb]{0.9,0.9,0.9}{
            \makecell[{{p{\linewidth}}}]{
                \texttt{\tiny{[GM$|$GM]}}
                \texttt{No} \\
            }
        }
    }
    & & \\ \\

    \theutterance \stepcounter{utterance}  
    & & & \multicolumn{2}{p{0.3\linewidth}}{
        \cellcolor[rgb]{0.9,0.9,0.9}{
            \makecell[{{p{\linewidth}}}]{
                \texttt{\tiny{[GM$|$GM]}}
                \texttt{no} \\
            }
        }
    }
    & & \\ \\

    \theutterance \stepcounter{utterance}  
    & & & \multicolumn{2}{p{0.3\linewidth}}{
        \cellcolor[rgb]{0.9,0.9,0.9}{
            \makecell[{{p{\linewidth}}}]{
                \texttt{\tiny{[GM$|$GM]}}
                \texttt{game\_result = WIN} \\
            }
        }
    }
    & & \\ \\

\end{supertabular}
}

\end{document}
